%#! latex main.tex
%\appendix
\chapter{Overview of theory of $p$-adic uniformization.}
\renewcommand{\thechapter}{\Alph{chapter}}
\addtocontents{toc}{\protect\par\vskip2mm\hskip80mm by F. Kato \par\vskip5mm}

\vspace{1ex}
\begin{center}
By {\sc F.\ Kato}
\end{center}

\vspace{2ex}
In this appendix we always denote by $K$, $R$, and $\pi$ a complete 
discrete valuation field, its valuation ring, and a prime element of
$R$, respectively. 
We assume that the residue field $k=R/\pi R$ is finite and consists of 
$q$ elements.
The Mumford-Kurihara-Mustafin uniformization is a procedure to
construct nice analytic (and in many cases algebraic) varieties by
taking discrete Schottky-type quotients of a certain $p$-adic analogue
of symmetric space, so-called {\it Drinfeld symmetric space} $\Omega$, 
or its variants.

%
\section{Bruhat-Tits building.}
%

First we survey the construction of a certain simplicial complex, called 
Bruhat-Tits building (attached to $\mathrm{PGL}(n+1,K)$), which will be 
closely related with the Drinfeld symmetric space.
Let $n$ be a positive integer and $V$ an $n+1$ dimensional vector space
over $K$.
A {\it lattice} in $V$ is a finitely generated 
$R$-submodule of $V$ which spans $V$ over $K$. 
Every lattice is therefore a free $R$-module of rank $n+1$.
Let $\til{\Delta}_0$ be the set of all lattices in $V$. 
We say that two lattices $M_1$ and $M_2$ are {\it similar} if there
exists $\lambda\in K^{\times}$ such that $M_1=\lambda M_2$.
The similarity is obviously an equivalence relation.
We denote by
$\Delta_0$ the set of all similarity classes of lattices in $V$.

\begin{appdfn}
The {\it Bruhat-Tits building} (attached to $\mathrm{PGL}(n+1,K)$) is
the finite dimensional simplicial complex $\Delta$ with the vertex set
$\Delta_0$ defined as follows: A finite subset $\{\Lambda_0,\ldots,
\Lambda_l\}$ of $\Delta_0$ forms an $l$-simplex if and only if, after 
permuting indices if necessary, one can choose $M_i\in\Lambda_i$ for 
$0\leq i\leq n+1$ such that
$$
M_0\supsetneq M_1\supsetneq\cdots\supsetneq M_l\supsetneq\pi M_0.
$$
\end{appdfn}

The role of $M_0$ is by no means important; for example, one can shift
the indices like
$M_1\supsetneq\cdots\supsetneq M_l\supsetneq\pi M_0\supsetneq\pi M_1$.

To understand the structure of $\Delta$, let us fix one vertex
$\Lambda=[M]$. 
Suppose $\{\Lambda_0,\ldots,\Lambda_l\}$ is an $l$-simplex having
$\Lambda$ as a vertex; we may assume $\Lambda=\Lambda_0$, and can take
$M_i\in\Lambda_i$ for $1\leq i\leq l$ such that 
$M\supsetneq M_1\supsetneq\cdots\supsetneq M_l\supsetneq\pi M$.
Set $\ovl{M}_i\colon=M_i/\pi M$.
Then we get a flag 
$$
\ovl{M}\supsetneq\ovl{M}_1\supsetneq\cdots\supsetneq\ovl{M}_l\supsetneq 0
$$
of length $l+1$ consisting of subspaces of the $n+1$-dimensional
$k$-vector space $\ovl{M}=M/\pi M$.
Moreover, by the elementary algebra one finds that this yields a
one-to-one correspondence between $l$-simplices containing $\Lambda$ and 
length $l+1$ frags in $\ovl{M}=M/\pi M$.
By this quite easy observation one knows
\begin{enumerate}
\item $\Delta$ is an $n$-dimensional locally finite simplicial complex,
\item every simplex in $\Delta$ is a face of a chamber, a simplex having 
      maximal ($=n$) dimension.
\end{enumerate}
Furthermore, by the construction,
\begin{enumerate}
\setcounter{enumi}{2}
\item the group $\Aut_K(V)/K^{\times}(\cong\mathrm{PGL}(n+1,K))$
      naturally acts on $\Delta$ simplicially.
\end{enumerate}

\begin{appexa}[$n=1$ case]
In $n=1$ case, the vertices adjacent to a fix vertex $\Lambda$ are in 
one-to-one correspondence with the set of all $k$-rational points of
the projective line $\mathbb{P}^1_k$; in particular, each vertex has 
exactly $q+1$ adjacent vertices.
Moreover, it can be seen that the simplicial complex $\Delta=\Delta^1$
is a {\it tree}.
The following picture shows how it looks like in case $q=2$:
$$
\setlength{\unitlength}{1pt}
\begin{picture}(200,100)(0,0)
\put(100,50){\circle*{3}}
\put(100,50){\line(0,1){20}}\put(100,70){\circle*{3}}
\put(100,50){\line(2,-1){20}}\put(120,40){\circle*{3}}
\put(100,50){\line(-2,-1){20}}\put(80,40){\circle*{3}}
%
\put(100,70){\line(2,1){12}}\put(112,76){\circle*{3}}
\put(100,70){\line(-2,1){12}}\put(88,76){\circle*{3}}
\put(120,40){\line(0,-1){12}}\put(120,28){\circle*{3}}
\put(120,40){\line(2,1){12}}\put(132,46){\circle*{3}}
\put(80,40){\line(0,-1){12}}\put(80,28){\circle*{3}}
\put(80,40){\line(-2,1){12}}\put(68,46){\circle*{3}}
%
\multiput(112,76)(0,4){3}{\circle*{.3}}
\multiput(112,76)(4,-2){3}{\circle*{.3}}
\multiput(88,76)(0,4){3}{\circle*{.3}}
\multiput(88,76)(-4,-2){3}{\circle*{.3}}
\multiput(120,28)(4,-2){3}{\circle*{.3}}
\multiput(120,28)(-4,-2){3}{\circle*{.3}}
\multiput(132,46)(0,4){3}{\circle*{.3}}
\multiput(132,46)(4,-2){3}{\circle*{.3}}
\multiput(80,28)(4,-2){3}{\circle*{.3}}
\multiput(80,28)(-4,-2){3}{\circle*{.3}}
\multiput(68,46)(0,4){3}{\circle*{.3}}
\multiput(68,46)(-4,-2){3}{\circle*{.3}}
\end{picture}\vspace{-1ex}
$$
\end{appexa}

In general, the geometrical realization $|\Delta|$ of $\Delta$ is known
to be a contractible topological space.

\begin{appexa}[$n=2$ case]\label{appexa-B}
Already in case $n=2$, the structure of $\Delta=\Delta^2$ is fairly
complicated.
For a fixed vertex $\Lambda$, the set of all vertices adjacent of 
$\Lambda$ is divided in two; one for those corresponding to $q^2+q+1$
$k$-rational points in $\mathbb{P}^2_k$ and the other for those
corresponding to $q^2+q+1$ $k$-rational lines in $\mathbb{P}^2_k$.
So it is convenient to consider the surface $B$ which is obtained by
blowing-up of $\mathbb{P}^2_k$ at every $k$-rational points.
Two vertices adjacent to $\Lambda$ forms with $\Lambda$ a $2$-simplex
if and only if the corresponding lines in $B$ intersect.
\end{appexa}

%
\section{Drinfeld symmetric space.} 
%

To each lattice $M$ in $V$ we associate an $R$-scheme
$$
\mathbb{P}(M)\colon=\mathrm{Proj}(\mathrm{Sym}_RM).
$$
Obviously the scheme $\mathbb{P}(M)$ depends only on the similarity
class; so we should write as $\mathbb{P}(\Lambda)$, where $\Lambda=[M]$.
The scheme $\mathbb{P}(M)$ is isomorphic to $\mathbb{P}^n_R$; but the
isomorphism cannot be taken canonically, while there exists the
{\it canonical} isomorphism $\mathbb{V}\cong\mathbb{P}(M)\otimes_RK$ 
due to the canonical isomorphism $M\otimes_RK\cong V$.
So we get a bijection:
$$
\Delta_0\stackrel{\sim}{\longrightarrow}\left\{
\begin{array}{c|l}
(\mathbb{P},\phi) &\mathbb{P}\stackrel{\sim}{\rightarrow}\mathbb{P}^n_R,\\
&\phi\colon\mathbb{P}\otimes_RK\stackrel{\sim}{\rightarrow}\mathbb{P}(V)
\end{array}
\right\}/\sim,
$$
where $(\mathbb{P}_1,\phi_1)\sim(\mathbb{P}_2,\phi_2)$ if and only if
there exists an $R$-isomorphism
$\Phi\colon\mathbb{P}_1\stackrel{\sim}{\rightarrow}\mathbb{P}_2$
which makes the diagram 
$$
\begin{array}{ccc}
\mathbb{P}_1\otimes_RK&\stackrel{\Phi\otimes_R K}
{\longrightarrow}&\mathbb{P}_2\otimes_RK\\
\llap{$\phi_1$}\downarrow&&\downarrow\rlap{$\phi_2$}\\
\mathbb{P}(V)&=&\mathbb{P}(V)
\end{array}
$$
commutative.

Next we consider two vertices $\Lambda_1$ and $\Lambda_2$ which are 
adjacent with each other.
The canonical isomorphism between the generic fibers induce a natural
birational map between $\mathbb{P}(\Lambda_1)$ and 
$\mathbb{P}(\Lambda_2)$.
More explicitly, taking representatives $M_i\in\Lambda_i$ ($i=1,2$) so
that $M_1\supsetneq M_2\supsetneq\pi M_1$, we have the birational map
from $\mathbb{P}(\Lambda_1)$ to $\mathbb{P}(\Lambda_2)$ induced by
$M_2\hookrightarrow M_1$.
We define the {\it join}
$\mathbb{P}(\Lambda_1)\vee\mathbb{P}(\Lambda_2)$ to be the closure of
the graph in $\mathbb{P}(\Lambda_1)\times_R\mathbb{P}(\Lambda_2)$ of
this rational map.
It is easily seen that the join 
$\mathbb{P}(\Lambda_1)\vee\mathbb{P}(\Lambda_2)$ is $R$-isomorphic to
the blow-up of $\mathbb{P}(\Lambda_1)$ at the closed subscheme 
$\mathrm{Proj}(\mathrm{Sym}_kM_1/M_2)$.
In particular, the closed fiber of 
$\mathbb{P}(\Lambda_1)\vee\mathbb{P}(\Lambda_2)$ consists of two
rational varieties intersecting transversally.

\begin{applem}
The operation $\vee$ is associative and commutative.
\end{applem}

The proof is straightforward, and left to the reader. 
By the lemma, we can extend the construction to any finite subcomplex of
$\Delta$.
More precisely, for a finite subcomplex $\matheur{S}$, we take the joins 
of all vertices in $\matheur{S}$; the resulting $R$-scheme is denoted by 
$\mathbb{P}(\matheur{S})$.
The scheme $\mathbb{P}(\matheur{S})$ is a regular projective scheme over
$R$ having the following properties:
\begin{enumerate}
\item The generic fiber is canonically isomorphic to $\mathbb{P}(V)$.
\item The closed fiber is a reduced normal crossing of non-singular
      rational varieties of which the dual graph coincides with
      $\matheur{S}$.
\end{enumerate}

We want to do this for a general convex subcomplex of $\Delta$; this cannot 
be done in terms of schemes because we have to consider a limit of
blow-up's of $\mathbb{P}^n_R$ centered in the closed fiber.
But we can do this in terms of formal schemes, or rigid analytic spaces.
Let $\Delta_{\ast}$ be a convex subcomplex of $\Delta$.
For any finite subcomplex $\matheur{S}$ of $\Delta_{\ast}$ we define 
$\widehat{\Omega}(\matheur{S})$ to be the completion of
$\mathbb{P}(\matheur{S})$ along the closed fiber.
Consider the maximal Zariski open subset $\widehat{\Omega}(\matheur{S})'$
of $\widehat{\Omega}(\matheur{S})$ such that, for any finite subcomplex
$\matheur{T}$ of $\Delta_{\ast}$ containing $\matheur{S}$, the induced
morphism $\rho^{\matheur{T}}_{\matheur{S}}\colon\widehat{\Omega}(\matheur{T})
\rightarrow\widehat{\Omega}(\matheur{S})$ gives an isomorphism restricted to 
the pull-back of $\widehat{\Omega}(\matheur{S})'$.

\begin{appdfn}
For a convex subcomplex $\Delta_{\ast}$ of $\Delta$ we define
$$
\widehat{\Omega}(\Delta_{\ast})\colon=\bigcup_{\matheur{S}\colon
\textrm{\scriptsize{finite}}}
\widehat{\Omega}(\matheur{S})',
$$
where $\matheur{S}$ runs through all finite subcomplex of
$\Delta_{\ast}$.
The corresponding rigid analytic space (Raynaud generic fiber) is
denoted by $\Omega(\Delta_{\ast})$.
In case $\Delta_{\ast}=\Delta$ we simply write $\widehat{\Omega}=
\widehat{\Omega}(\Delta)$ and $\Omega=\Omega(\Delta)$.
The rigid analytic space $\Omega$ is called the {\it Drinfeld symmetric
space}.
\end{appdfn}

The formal scheme $\widehat{\Omega}(\Delta_{\ast})$ is formally locally of
finite type, of which the closed fiber is the normal crossings of
non-singular rational varieties having the dual graph isomorphic to 
$\Delta_{\ast}$.
By definition, the group $\mathrm{PGL}(n+1,K)$ naturally acts on 
$\widehat{\Omega}$, and hence $\Omega$.

\begin{appexa}
Let us consider $\widehat{\Omega}$ in $n=2$. 
The closed fiber $\widehat{\Omega}_0$ consists of countably many 
components which are all isomorphic to the surface $B$ considered 
in \ref{appexa-B}. 
In each component every $k$-rational line is a double curve in 
$\widehat{\Omega}_0$, and every $k$-rational point is a triple point.
\end{appexa}

{\it Structure of $\Omega$.}\ 
First we note that, as a point set, we have an equality
$$
\Omega=\mathbb{P}^{n,\mathrm{an}}_K\setminus\bigcup_{\matheur{H}\in
\mathcal{H}}\matheur{H},
$$
where $\mathcal{H}$ is the set of all $K$-rational hyperplanes.

Here we do not attempt to give a proof of this fact, but rather give a 
sketchy explanation in order to get a feeling.
We limit ourselves to $n=1$ case just for simplicity.
The formal scheme $\widehat{\Omega}$ is, roughly speaking, the {\it limit}
of successive blow-up's 
$$
\cdots\longrightarrow U_i\stackrel{\rho_i}{\longrightarrow}U_{i-1}
\longrightarrow\cdots\longrightarrow U_0
$$
with $U_0=\mathbb{P}^1_R$ and $\rho_i$ the blow-up centered at all the 
$k$-rational points in the closed fiber of $U_{i-1}$.
Let $z$ be a point in the generic fiber $\mathbb{P}^1_K$ of each $U_i$.
We consider the closure $\ovl{z}$ of $z$ in each $U_i$.
Since $U_i$ is a regular scheme, the section $\ovl{z}$ intersects the
closed fiber at a smooth point transversally. 
If $z$ is {\it not} a $K$-rational point, then for sufficiently large 
$N$, $\ovl{z}$ intersects the closed fiber of $U_i$ at a non-$k$-rational
smooth point for $i>N$.
But, if $z$ is $K$-rational, $\ovl{z}$ intersects the closed
fiber always at a $k$-rational point which is in the next blow-up center.
So if we consider the limit, such formal sections determined by $K$-rational 
points does not exist any more, while those by non-$K$-rational points 
certainly exist.
This explains the above equality in $n=1$.
For higher $n$ we can apply the similar argument.

Next we discuss the analytic structure of $\Omega$; viz.\ we show how
to introduce an admissible covering to $\Omega$; since we have
constructed the rigid analytic space $\Omega$ from the formal model, we 
must have one such covering such that the associated formal model 
recovers $\widehat{\Omega}$.
This can be actually described in an elegent way in terms of norms: 
A non-negative real valued function $\alpha\colon V\rightarrow
\mathbb{R}_{\geq 0}$ is said to be a {\it norm} over $K$ if it
satisfies:
\begin{enumerate}
\item $\alpha(x)>0$ whenever $x\neq 0$.
\item For $a\in K$, $\alpha(ax)=|a|\alpha(x)$.
\item $\alpha(x+y)\leq\mathrm{max}\{\alpha(x),\alpha(y)\}$.
\end{enumerate}
Two norms $\alpha$ and $\alpha'$ over $K$ is said to be {\it similar} if 
$\alpha'=\lambda\alpha$ for some $\lambda\in\mathbb{R}_{>0}$.
The set of all similarity classes of norms on $V$ over $K$ forms in an
obvious way a topological space. 
The following fact is well-known (\cite{tsopn}):

\begin{apppro}
The topological space of similarlity classes of norms on $V$ over $K$ is
$\mathrm{PGL}(n+1,K)$-equivaliantly homeomorphic to the topological 
realization $|\Delta|$ of the Bruhat-Tits building $\Delta$.
\end{apppro}

The homeomorphism is given as follows: To any lattice $M$ we need to 
associate a norm $\alpha_M$.
For a non-zero $x\in V$, the set $\{a\in K\ |\ ax\in M\}$ is a
fractional ideal of $R$, and hence is of form $(\pi^m)$.
We then define $\alpha_M(x)=q^m$ for non-zero $x\in V$, and
$\alpha_M(0)=0$.
It is easily verified that $\alpha_M$ is a norm.
The similarity class of $\alpha_M$ depends only on the similarity class
of $M$.
Thus we get a mapping from $\Delta_0$ to the set of similarity classes
of norms.
Let $\Lambda_0$ and $\Lambda_1$ be vertices adjacent with each other.
We can choose a basis $e_0,\ldots,e_n$ of $V$ such that 
$M_0=\bigoplus^n_{i=0}Re_i$ and $M_1=\bigoplus^{n-1}_{i=0}Re_i\oplus
R\pi e_n$ give representatives of $\Lambda_0$ and $\Lambda_1$,
respectively.
Then the norms $\alpha_{M_0}$ and $\alpha_{M_1}$ are given by
\begin{eqnarray*}
\alpha_{M_0}(\sum^n_{i=0}a_ie_i)&=&\mathrm{max}\{|a_0|,\ldots,|a_n|\},\\
\alpha_{M_1}(\sum^n_{i=0}a_ie_i)&=&\mathrm{max}\{|a_0|,\ldots,|a_{n-1}|,
|\pi|^{-1}|a_n|\},
\end{eqnarray*}
respectively.
For $0<t<1$, the class of the norm
$$
\alpha_t(\sum^n_{i=0}a_ie_i)=\mathrm{max}\{|a_0|,\ldots,|a_{n-1}|,
|\pi|^{-t}|a_n|\}
$$
corresponds to the point $t\Lambda_0+(1-t)\Lambda_1$ on the edge
connecting $\Lambda_0$ and $\Lambda_1$.

Now, to any point $z=(z_0\colon\cdots\colon z_n)\in\mathbb{P}(V)$
(subject to some $K$-basis $\{e_0,\ldots,e_n\}$of $V$) 
we associate, up to invertible factor, a non-negative real
valued function
$$
\alpha_z(\sum^n_{i=0}a_ie_i)=\biggm|
\sum_{i=0}^n a_iz_i\biggm|
$$
on $V$.
Note that 
{\it this determines a class of norm if and only if $z$ is not lying in any 
$K$-rational hyperplane.}
Thus we get a mapping 
$$
\rho\colon\Omega\longrightarrow |\Delta|.
$$
The admissible covering in question is actually given by the collection
of subsets of $\Omega$ consisting of pull-backs of each simplices of
$|\Delta|$ by $\rho$.
For example, notation being as above, 
$$
\rho^{-1}(\Lambda_0)=\{|z_0|=\cdots=|z_n|=1\}\setminus
\bigcup_{H}\{(z_0\colon\cdots\colon z_n)\ \mathrm{mod}\ \pi\in H\},
$$
where $H$ runs through all hyperplanes in $\mathbb{P}^n_k$, and 
$$
\rho^{-1}(t\Lambda_0+(1-t)\Lambda_1)=\{|z_0|=\cdots=|z_{n-1}|=1,
|z_n|=|\pi|^{-t}\}
$$
for $0<t<1$.

%
\section{Uniformization.}
%

Recall that, once we fix a $K$-basis of $V$, the group 
$\mathrm{PGL}(n+1,K)$ canonically acts on the Bruhat-Tits building 
$\Delta$ and the Drinfeld symmetric space $\Omega$.
We shall consider discrete groups in $\mathrm{PGL}(n+1,K)$ and the
quotients of $\Delta$ and $\Omega$ by them.
A subgroup $\Gamma\in\mathrm{PGL}(n+1,K)$ is said to be {\it
hyperbolic} if it acts on $\Delta$ discretely and freely.

To a hyperbolic subgroup $\Gamma$ we associate a convex subcomplex 
$\Delta_{\Gamma}$ of $\Delta$ by the following recipe:
In general, a convex subcomplex generated by vertices of form
$$
\left[\sum^n_{i=0}R\cdot\pi^{\sigma_i}\matheur{Y}_i\right],
\qquad
\sigma_i\in\Z,\qquad 0\leq i\leq n,
$$
subject to a basis $\matheur{Y}_0,\ldots,\matheur{Y}_n$ of $V$, is said
to be an {\it apartment}.
The geometrical realization of each apartment is a triangulation of 
$\mathbb{R}^n$ of $\mathrm{A}_n$-type, and $\Delta$ is the union of all
apartments.
There exists an obvious bijection between the set of all apartments and
the set of all $K$-split maximal tori in $\mathrm{PGL}(n+1,K)$.
For any $K$-split maximal tori $T$ the untersection $\Gamma\cap T$ is a
commutative discrete group of rank at most $n$.
If this rank is $n$ we call the apartment corresponding to $T$ a 
{\it $\Gamma$-apartment}, and we define $\Delta_{\Gamma}$ to be the
union of all $\Gamma$-apartments.
Clearly the group $\Gamma$ acts on $\Delta_{\Gamma}$; if the quotient
$\Delta_{\Gamma}/\Gamma$ is a finite simplicial complex, we say 
$\Gamma$ is {\it normal hyperbolic}.

Now let us be given a normal hyperbolic subgroup $\Gamma$.
For each $\gamma\in\Gamma$ and each finite subcomplex $\matheur{S}$ of
$\Delta_{\Gamma}$, $\gamma\matheur{S}$ is again a finite subcomplex of
$\Delta_{\Gamma}$, and then, $\gamma$ induces an $R$-isomorphism 
$\mathbb{P}(\matheur{S})\cong\mathbb{P}(\gamma\matheur{S})$ which
induces on generic fibers ($=\mathbb{P}^n_K$) exactly the linear 
automorphism $\gamma$.
Considering all $\matheur{S}$, gluing, and taking Raynaud generic fiber
give a rigid analytic automorphism of $\Omega(\Delta_{\Gamma})$ which 
we denote again by $\gamma$.
Thus we get an homomorphism 
$$
\Gamma\longrightarrow\mathrm{Aut}_{K\textrm{\scriptsize{-rig.}}}
\Omega(\Delta_{\Gamma}).
$$

\begin{appthm}[Mustafin\ \cite{nu}, Kurihara\ \cite{copubathp}]
Let $\Gamma$ be a normal hyperbolic subgroup of $\mathrm{PGL}(n+1,K)$
and $\Delta_{\Gamma}$ as above.
Then there exists, unique up to isomorphisms,
a rigid analytic space $X_{\Gamma}$, smooth, proper and of finite type
over $K$, and a topological covering map $\matheur{p}\colon
\Omega(\Delta_{\Gamma})\rightarrow X_{\Gamma}$ such that
\begin{enumerate}
\item $\matheur{p}\circ\gamma=\matheur{p}$ for any $\gamma\in\Gamma$,
\item for $x,y\in\Omega(\Delta_{\Gamma})$, $\matheur{p}(x)=\matheur{p}
(y)$ if and only if $x=\gamma(y)$ for some $\gamma\in\Gamma$.
\end{enumerate}
\end{appthm}

In the references, this theorem has been stated in terms of formal
schemes; viz.\ they took quotients of
$\widehat{\Omega}(\Delta_{\Gamma})$ to get a formal scheme 
$\widehat{X}_{\Gamma}$ flat and formally of finite type over $R$.
We have an admissible covering of $X_{\Gamma}$ induced from that of 
$\Omega(\Delta_{\Gamma})$; in terms of formal scheme, the closed fiber
of $\widehat{X}_{\Gamma}$ is a divisor of normal crossings with rational
components having the dual group isomorphic to $\Delta_{\Gamma}/\Gamma$.

\vspace{1ex}
{\it Algebraicity.}\ In $n=1$ case the rigid analytic curve $X_{\Gamma}$
is algebraizable for any $\Gamma$ (\cite{aacodcoclr}).
In higher dimensions, we know the following sufficient condition:
Let $\Gamma$ be a torsion-free {\it uniform lattice} in 
$\mathrm{PGL}(n+1,K)$, viz.\ a
discrete co-compact subgroup of finite co-volume without elements of 
finite order.
Then, since it is finite co-volume, we have $\Delta_{\Gamma}=\Delta$.
Hence in this case, we are going to take a quotient of the Drinfeld
symmetric space $\Omega$.

\begin{appthm}[loc.\ cit.]
If $\Gamma$ is a torsion-free uniform lattice, then the quotient rigid
analytic variety $X_{\Gamma}$ can be algebraizable to a non-singular
projective variety having the ample canonical class.
\end{appthm}

In terms of formal scheme, this means that the relative canonical sheaf
of $\widehat{X}_{\Gamma}/\mathrm{Spf}R$ is ample.

%
\section{Examples.}

{\it Tate curve.}\ 
Let $T$ be the maxial torus consisting of all invertible diagonal
matrices. We take $n$ elements $q_i\in K^{\times}$ ($1\leq i\leq n$)
with $|q_i|<1$ and set $q_0=1$.
Define $\Gamma\subset T$ to be the set of all diagonal matrices with
$(i,i)$-entry an integer power of $q_i$ for $0\leq i\leq n$.
Then $\Gamma$ is a normal hyperbolic subgroup of $\mathrm{PGL}(n+1,K)$
with $\Delta_{\Gamma}$ the apartment corresponding to $T$.
If $n=1$, the curve $\Omega(\Delta_{\Gamma})/\Gamma$ is known to be 
the so-called {\it Tate curve}; in this case
$\Omega(\Delta_{\Gamma})$ is simply $\mathbb{P}^1_K$ minus two points.

\vspace{1ex}
{\it Mumford curve.}\ 
Consider again the $n=1$ case.
The rigid analytic curve $X_{\Gamma}$ with $\Gamma$ generated by at
least two elements is called a {\it Mumford curve}.
A Mumford curve is, in fact, algebraized to a non-singular projective
curve of genus greater than $1$.
Moreover, the formal model provided by the admissible covering inhereted 
from $\Omega(\Delta_{\Gamma})$ gives a stable $k$-split multiplicative
reduction of that curve.
The converse is also known:

\begin{appthm}[Mumford\ \cite{aacodcoclr}]
Let $X$ be a smooth and proper rigid analytic curve over $K$, and
suppose $X$ admits a formal model with $k$-split multiplicative
reduction. Then $X$ is isomorphic to $X_{\Gamma}$ for a unique 
$\Gamma$.
\end{appthm}

{\it Surface case.}\ 
In higher dimension it becomes extremely difficult to find normal
hyperbolic groups; but there are several known examples of torsion-free
uniform lattices.
Let $\Gamma$ is a torsion-free uniform lattice, and $N$ the number of
$\Gamma$-orbits in the vertex set $\Delta_0$.
Mumford \cite{aaswkak9pq0} calculated the several numerical invariants of the
quotient surface $X_{\Gamma}$:
\begin{enumerate}
\item $\chi(\OO_{X_{\Gamma}})=N(q-1)^2(q+1)/3$.
\item $(\matheur{c}_{1,X_{\Gamma}})^2=3\matheur{c}_{2,X_{\Gamma}}=
3N(q-1)^2(q+1)$.
\item $\matheur{q}(X_{\Gamma})=0$,
\end{enumerate}
where $\matheur{q}$ denotes the irregularity.
It is very interesting that the Mumford-Kurihara-Mustafin uniformization 
in two dimension always produces the surfaces which satisfies
the equality in Miyaoka-Yau inequality.
Mumford \cite{aaswkak9pq0} constructed one example of $\Gamma$ in $q=2$ such
that $N=1$; the resulting surface is one of the so-called ``fake
projective planes''.
Recently, Cartwright, Mantero, Steger, and Zappa constructed several
examples of lattices in $q=2$ and $q=3$. 
Making use of this groups, in \cite{tsrtfnu}, Ishida and the author
discussed other possible fake projective planes, and showed that there
are at least three fake projective planes (incl.\ Mumford's one) which
are not isomorphic to among others.
We will see in more detail the construction of Mumford's fake projective
plane in the next section.

\vspace{1ex}
{\it Shimura variety case.}
Let $D$ be a quaternion algebra over a totally real field $F$.
We assume that $D$ ramifies at every infinite place except one $\infty_0$.
Let $\Gamma_{\infty}$ be a subgroup of $D^{\times}/F^{\times}$ defined
by some congruence condition; we shall assume that it is ``sufficiently
small''.
Then it is classically known that $\Gamma_{\infty}$ acts discontinuously 
on the upper $\frac{1}{2}$ plane $\matheur{H}$ such that the quotient
$\matheur{H}/\Gamma_{\infty}$ is a {\it projective} curve (Shimura curve
of PEL-type).
The canonical model $\matheur{Sh}$ is defined over some ray-class field
of $F$.

There must be a finite place $\nu$ at which $B$ ramifies.
Let $p$ be the residue characteristic of $\nu$.
We assume that $\Gamma_{\infty}$ is maximal at $\nu$.
Then the Shimura curve $\matheur{Sh}$ has a bad reduction 
at $\nu$, and due to Cherednik \cite{uoacbdsopwcq} and Drinfeld
\cite{copsr}, the associated rigid analytic space over 
$F^{\mathrm{ur}}_{\nu}$ has a $p$-adic uniformization, where
$F^{\mathrm{ur}}_{\nu}$ is the maximal unramified extension of
$F_{\nu}$.
Let $\ovl{D}$ be the quaternion algebra which has local invariants
obtained precisely by switching those of $D$ at $\nu$ and $\infty_0$
Then there exists a subgroup $\Gamma_{\nu}$ in 
$\ovl{D}^{\times}/F^{\times}$ defined basically by the same congruence 
relation as that of $\Gamma_{\infty}$ such that
$\Omega\otimes F^{\mathrm{ur}}/\Gamma_{\nu}$ is isomorphic to the
assocaied analytic space of $\matheur{Sh}$ at $\nu$.
(This isomorphism actually decends to some finite unramified extension
of $F_{\nu}$, depending on the data of connected components of
$\matheur{Sh}$.)

This story has been generalized to higher dimension by Rapoport-Zink
\cite{psfpg} and Boutot-Zink \cite{tpuosc}.
The next section is devoted to explain one special example of the 
$p$-adic uniformization of a Shimura surface.

%
\section{Mumford's fake projective plane.}
%
Our example of $p$-adically uniformizable Shimura surface is constructed 
by means of some division algebra $D$. Th following construction is
inspired by the Mumford's oroginal paper \cite{aaswkak9pq0}.

Let $\zeta=\zeta_7$ be a primitive $7$th-root of unity, and
set $L\colon=\Q(\zeta)$; $L$ is a cyclic extension of $\Q$ of degree $6$
having an intermediate quadratic extension $K\colon=\Q(\lambda)\ (\cong
\Q(\sqrt{-7}))$, where $\lambda=\zeta+\zeta^2+\zeta^4$.
The Galois group of $L/K$ is generated by the Frobenius map $\sigma\colon
\zeta\mapsto\zeta^2$, and that of $K/\Q$ is generated by the complex 
conjugation $z\mapsto\ovl{z}$.
Note that the prime $2$ decomposes on $K$, more explicitly, 
$2=\lambda\ovl{\lambda}$ gives the prime factorization, and the prime
$7$ ramifies.
We fix an infinite place $\varepsilon\colon K\hookrightarrow\C$ by
$\lambda\mapsto\frac{-1+\sqrt{-7}}{2}$.

We set $\mu\colon=\lambda/\ovl{\lambda}$, and define the central
division algebra $D$ over $K$ of dimension $9$ by
$$
D\colon=\bigoplus^2_{i=0}L\Pi^i
$$
with $\Pi^3=\mu$ and $\Pi z=z^{\sigma}\Pi$ for $z\in L$.
The involution $\ast$ on $D$ is defined by $z^{\ast}=\ovl{z}$ for $z\in L$
and $\Pi^{\ast}=\ovl{\mu}\Pi^2$.
It is easily verified that $\ast$ is positive.
%We set
%$$
%\mathcal{O}_D\colon=\mathcal{O}_L\oplus\mathcal{O}_L\ovl{\lambda}\Pi\oplus
%\mathcal{O}_L\ovl{\lambda}\Pi^2,
%$$
%which is clealy stable under $\ast$.
%Moreover, one can easily verify that $\mathcal{O}_{D_{\lambda}}$ and
%$\mathcal{O}_{D_{\ovl{\lambda}}}$ are maximal.

We introduce a non-degenerate anti-symmetric $\Q$-bilinear
form $\psi=\psi_b$ on $V=D$ by putting 
$b=(\lambda-\ovl{\lambda})-\ovl{\lambda}\Pi+\ovl{\lambda}\Pi^2$.
Then define an algebraic group $G$ over $\Q$ as follows: 
Given a $\Q$-algebra $R$, we set
$$
G(R)\colon=\{g\in((D\otimes R)^{\mathrm{opp}})^{\times}\ |\ \psi(xg,yg)=
c(g)\psi(x,y),\ c(g)\in R^{\times},\ x,y\in V\otimes R\}.
$$
The algebraic group $G_{\R}$ is isomorphic to $\mathrm{GU}(2,1)$, which 
acts on the two dimensional complex unit-ball $\matheur{B}$.
In the following we will construct an open compact subgroup $C$ in
$G(\mathbb{A}_{\mathrm{f}})$.
By these data, we can define a Shimura variety
$$
\mathcal{S}_C\colon=G(\Q)\setminus\matheur{B}\times 
G(\mathbb{A}_{\mathrm{f}})/C.
$$
The canonical model of $\mathcal{S}_C$ is defined over $E\colon=
\varepsilon(K)$, which we denote by $\matheur{Sh}_C$.

\vspace{1ex}
{\it The congruence condition.}\ 
First we define an algebraic group $I$ over $\Q$ by 
$$
I(R)\colon=\{g\in(\mathrm{M}_d(K)\otimes R)^{\times}\ |\
H^{-1}\,^t\ovl{g}H\in R^{\times}\}
$$
for any $\Q$-algebra $R$, where 
$$
H=\left[
\begin{array}{lll}
3&\ovl{\lambda}&\ovl{\lambda}\\
\lambda&3&\ovl{\lambda}\\
\lambda&\lambda&3
\end{array}
\right].
$$
Then $I$ is an inner form of $G$ so that 
$G(\mathbb{A}^2_{\mathrm{f}})\cong I(\mathbb{A}^2_{\mathrm{f}})$.

To define the level structure, it suffices to define the 
prime-$2$-part $C^2$.
We define
$C^2$ in the form $C^2=C^{2,7}C_7$ with $C_7\in G(\Q_7)\cong I(\Q_7)$
and $C^{2,7}$ being the maximal one:
$$
C^{2,7}\colon=\{g\in G(\mathbb{A}^{2,7}_{\mathrm{f}})\ |\ \Gamma^{2,7}g
\subseteq\Gamma^{2,7}\},
$$
with $\Gamma^{2,7}=\Gamma\otimes_{\Z}\widehat{\Z}^{2,7}$, where
$\widehat{\Z}^{2,7}=\lim\Z/m\Z$ with the projective limit taken over $m$ 
such that $(2,m)=(7,m)=1$.

Now the component $C_7$ is defined as follows:
we have an isomorphism
$$
G(\Q_7)\stackrel{\sim}{\longrightarrow}\{g\in\mathrm{GL}_3(\til{\Q}_7)
\ |\ g^{\ast}Hg=c(g)H,\ c(g)\in\Q^{\times}_7\},
$$
where $\til{\Q}_7$ is the ramified quadratic extension of $\Q_7$, and
$\ast$ is the matrix transposition followed by the Galois action.
Let $\til{C}_7$ be the maximal open compact subgroup of $G(\Q_7)$ 
consisiting of $g$ which maps $\Gamma\otimes_{\Z}\Z_7$ to itself.
We consider modulo $\sqrt{-7}$ reduction of $\til{C}_7$, which is a
subgroup $G_0$ in $\mathrm{GL}_3(\mathbb{F}_7)$.
The matrix $H$ mod $\sqrt{-7}$ is of rank $1$, and has $2$ dimensional
null space $N_0$.
Restricting elements in $G_0$ to $N_0$, we obtain a homomorphism
$$
\pi\colon\til{C}_7\longrightarrow\mathrm{GL}_2(\mathbb{F}_7).
$$
The last group has $2^63^27$ elements. Let $S$ be a $2$-Sylow subgroup
of $\mathrm{GL}_2(\mathbb{F}_7)\cap\{g\ |\ \mathrm{det}\,g=\pm 1\}$.
Now we set $C_7\colon=\pi^{-1}(S)$.

\begin{appthm}\label{thm-uniformization}
%The canonical model $\matheur{Sh}_C$ of the Shimura variety obtained by%
%the data as above is
%a geometrically connected non-singular projective surface over $E=
%\Q(\sqrt{-7})$ such that the base change $\matheur{Sh}_C\otimes_E
%E_{\lambda}$ is isomorphic over $\Q_2\cong E_{\lambda}$ to the Mumford's
%fake projective plane \cite{aaswkak9pq0}.

The canonical model $\matheur{Sh}_C$ of the Shimura variety obtained by
the data as above has exactly three geometrically connected components defined
over $\Q(\zeta_7)$. Moreover, for any connected component $S$, its
base change $S\otimes_{\Q(\zeta_7)}\Q_2(\zeta_7)$ is isomorphic to the
Mumford's fake projective plane \cite{aaswkak9pq0} tensored by $\Q_2(\zeta_7)$.
\end{appthm}

More precisely, the associated rigid analytic space to 
$\matheur{Sh}_C\otimes_EE_{\lambda}$ is isomorphic to
$\Omega/\Gamma$ with
$$
\Gamma=I(\Q)\bigcap(J(\Q_2)\times C^2)\subset J(\Q_2)\cong
\mathrm{PGL}(3,\Q_2),
$$
where $I(\Q)$ is considered as a subgroup of 
$I(\mathbb{A}_{\mathrm{f}})\cong J(\Q_2)\times 
G(\mathbb{A}^2_{\mathrm{f}})$.
The proof of the theorem is in \cite{opuofpp}, where two more 
fake projective planes are
discussed as examples of Shimura varieties.

